\clearpage\fancyhead[L]{\small JUNO.}\hypertarget{native-375-lower}{}\label{native-375-lower}
\begin{center}\textbf{JUNO.}\par\vspace{4pt}\hrule\vspace{4pt}\fontsize{10}{12}\selectfont Monatliche Correspondenz. Band XXVIII. Seite 574\;..\;578. 1813. December.\par\vspace{4pt}\hrule\vspace{4pt}\end{center}


\begin{flushright}Göttingen, den 8. 9. und 16. Dec. 1813.\end{flushright}
\noindent --- --- --- Dem Beobachten ist das Wetter, wie gewöhnlich in dieser Jahrszeit, sehr ungünstig
gewesen. Meine Zeitbestimmung mache ich jetzt, da correspondirende Sonnenhöhen jetzt nie
gelingen, \SourceLeaf{376}indem es nur auf einzelne Stunden zuweilen sich etwas aufklärt, durch Sternazimuthe, die ich
(am Tage) mit dem Theodolithen beobachte. Ich vergleiche $\alpha$ Aurigae jetzt gegen
Sonnenuntergang mit einem 4000 Meter entfernten Thurm, dessen Azimuth im Sommer bestimmt
wurde, und es geht recht gut damit.

\Needspace{12\baselineskip}\par\indent Die Juno habe ich nur ein einzigesmal beobachten können, aber die Beobachtung ist an sich
sehr gut, und am Tage der Opposition selbst gemacht. Der Planet hatte die achte Grösse.
\begin{center}\fontsize{10}{12}\selectfont\setlength{\tabcolsep}{4pt}\renewcommand{\arraystretch}{1.15}\begin{tabular}{l|c|c|c}
1813 & \shortstack{Mittlere Zeit\\in Göttingen} & \shortstack{Gerade\\Aufsteigung $\JunoSymbol$} & \shortstack{Südliche\\Abweichung $\JunoSymbol$}\\\hline
Nov. 19 & $9^{\mathrm u}\;46^{\mathrm m}\;46^{\mathrm s}$ & $60^\circ\;30'\;35''0$ & $3^\circ\;5'\;9''3$\\
\end{tabular}\end{center}
\Needspace{7\baselineskip}\par\indent Die scheinbare Position des verglichenen Sterns wurde aus der \emph{Hist. Cél.} bestimmt
\begin{center}AR. $=60^\circ\;47'\;50''0$\quad Decl. $=3^\circ\;0'\;16''0$ südl.\end{center}

\Needspace{12\baselineskip}\par\indent Ich habe aus dieser Beobachtung ganz unbedenklich die Opposition abgeleitet, da schwerlich
anderswoher Beobachtungen eintreffen werden.
\begin{center}\emph{Achte Opposition der Juno:}\end{center}
\begin{center}1813 November 19. $18^{\mathrm u}\;11^{\mathrm m}\;25^{\mathrm s}$ Mittlere Zeit in Göttingen\end{center}
\begin{center}\fontsize{10}{12}\selectfont\setlength{\tabcolsep}{2pt}\renewcommand{\arraystretch}{1.15}\begin{tabular}{lcr}
Wahre Länge & $.\;.\;.\;.\;.\;.\;.\;.\;.\;.\;.\;.\;.$ & $57^\circ\;33'\;58''0$\\
Geocentrische Breite & $.\;.\;.\;.\;.\;.\;.\;.\;.\;.$ & $23\;18\;46.8\ \text{südl.}$\\
\end{tabular}\end{center}
\Needspace{6\baselineskip}\begin{center}\emph{Fortsetzung vom 9. December $10^{\mathrm u}$.}\end{center}
\par\indent Heute Abend habe ich, bei starkem Nebel und Mondschein, zwar noch nicht den Platz, wo ich
die Juno am 19. November beobachtet habe, wieder nachgesehen, aber die Juno selbst wieder
aufgesucht. \emph{Verhältnissmässig} schien sie heute \emph{heller} als am 19. November,
wenigstens war sie bedeutend heller, als der Stern achter Grösse, womit sie heute verglichen
wurde, da sie am 19. November dem verglichenen Sterne achter Grösse kaum gleich kam. Ich
ersuche Herrn \textsc{Nicolai}, die Beobachtungen zu reduciren, und mit den nachher folgenden
Elementen zu vergleichen. Die Beobachtung beruht zwar nur auf drei kümmerlichen
Vergleichungen mit einem Sterne der \emph{Hist. Cél.}, welcher am 20. Januar 1798 am zweiten
Faden $3^{\mathrm u}\;42^{\mathrm m}\;48^{\mathrm s}5$ beobachtet ist, kann aber doch zur
Controle dienen.
\begin{center}
1813 Decemb. 9. $0^{\mathrm u}\;22^{\mathrm m}\;2^{\mathrm s}$ St. Z., $\JunoSymbol$ folgt auf den Stern
$3^{\mathrm m}\;49^{\mathrm s}75$ Zeit $=57'\;26''25$ Bogen

$0^{\mathrm u}\;32^{\mathrm m}\;4^{\mathrm s}$ --- $\JunoSymbol$ ist südlicher $5'\;0''9$.
\end{center}
\noindent Ungefähr folgt daraus
\begin{center}AR. $=56^\circ\;56\tfrac12'$\quad Decl. $=3^\circ\;56\tfrac12'$ südl.\end{center}

\Needspace{15\baselineskip}\par\indent Die Elemente, auf die Oppositionen von 1810, 1811, 1812, 1813 gegründet, sind folgende (X):
\begin{center}\fontsize{9}{11}\selectfont\setlength{\tabcolsep}{2pt}\renewcommand{\arraystretch}{1.15}\begin{tabular}{lcr}
Epoche 1814, Meridian von Göttingen & $.\;.\;.\;.\;.\;.$ & $65^\circ\;19'\;23''8$\\
Tägliche tropische Bewegung & $.\;.\;.\;.\;.\;.\;.\;.\;.$ & $812''709$\\
Logarithm. der halben Axe & $.\;.\;.\;.\;.\;.\;.\;.\;.\;.$ & $0.4267631$\\
Perihelium (1814) & $.\;.\;.\;.\;.\;.\;.\;.\;.\;.\;.\;.\;.\;.$ & $53^\circ\;10'\;0''40$\\
Excentricitätswinkel & $.\;.\;.\;.\;.\;.\;.\;.\;.\;.\;.\;.\;.$ & $14\;43\;7.97$\\
Aufsteigender Knoten (1814) & $.\;.\;.\;.\;.\;.\;.\;.\;.\;.$ & $171\;10\;33.00$\\
Neigung & $.\;.\;.\;.\;.\;.\;.\;.\;.\;.\;.\;.\;.\;.\;.\;.\;.$ & $13\;4\;12.03$\\
\end{tabular}\end{center}
\par\indent Die Vergleichung dieser Elemente, besonders der mittlern Bewegung, mit den Elementen (VIII),
welche auf die Oppositionen von 1804, 1806, 1807, 1808 gegründet waren, zeigt, wie stark
sich die Jupiterstörungen äussern.

\Needspace{18\baselineskip}\begin{flushright}\leavevmode\SourceLeaf{377}Göttingen, den 16. December 1813.\end{flushright}
\noindent --- --- --- Ich habe mir doch die Mühe gegeben, auch der letzten IX. Opposition der Pallas die
Elemente anzupassen. Die Correctionen sind äusserst gering, nemlich:
\begin{center}\fontsize{10}{12}\selectfont\setlength{\tabcolsep}{2pt}\renewcommand{\arraystretch}{1.15}\begin{tabular}{lcr}
Länge des Knotens & $.\;.\;.\;.\;.\;.$ & $-2''7$\\
Neigung der Bahn & $.\;.\;.\;.\;.\;.$ & $-1.4$\\
Länge des Perihels & $.\;.\;.\;.\;.\;.$ & $-9.4$\\
Excentricitätswinkel & $.\;.\;.\;.\;.\;.$ & $+1.8$\\
Epoche 1810 & $.\;.\;.\;.\;.\;.\;.$ & $-4.2$\\
Tägliche Bewegung & $.\;.\;.\;.$ & $-0''00569$\\
Logarithm. der halben Axe & $.$ & $+0.0000021$\\
\end{tabular}\end{center}
\Needspace{22\baselineskip}\par\indent Die Uebereinstimmung der 9 Oppositionen mit den verbesserten Elementen ist nun folgende:
\begin{center}Unterschied:\end{center}
\begin{center}\fontsize{10}{12}\selectfont\setlength{\tabcolsep}{7pt}\renewcommand{\arraystretch}{1.15}\begin{tabular}{l|c|c}
Oppos. & \shortstack{Mittlere\\Länge} & \shortstack{Heliocentr.\\Breite.}\\\hline
I & $-1''3$ & $-5''3$\\
II & $-12.9$ & $-11.9$\\
III & $-20.4$ & $-0.4$\\
IV & $+2.5$ & $-6.8$\\
V & $+23.5$ & $+9.0$\\
VI & $+21.5$ & $-3.1$\\
VII & $+14.1$ & $-7.2$\\
VIII & $-16.1$ & $+4.0$\\
IX & $-10.9$ & $-0.0$\\
X & $-14.0$ & $+2.4$\\
\end{tabular}\end{center}
\par\indent Ich bin sehr neugierig, wie die Unterschiede ausfallen werden, wenn erst die
Saturnsstörungen dazu kommen. [Die Angabe für Oppos. X ist eine handschriftliche Bemerkung.
\textsc{Schering}.]
